\chapter{Deployment and Production}

\section*{Introduction}

This chapter describes the installation and configuration procedures for \caimp{}
in different usage contexts: local development, small-footprint VPS, and production
deployment.

\section{Deployment Modes}

\begin{table}[H]
\centering
\caption{Deployment modes and resource requirements}
\label{tab:deploy-modes}
\begin{tabularx}{\textwidth}{lXXX}
\toprule
\textbf{Mode} & \textbf{Command} & \textbf{Min. RAM} & \textbf{Use case} \\
\midrule
Full       & \texttt{make up}                      & 8\,GB & Production, development \\
Low-RAM    & \texttt{bash scripts/start-test.sh}   & 4\,GB & Economy VPS, CI \\
Production & Docker Swarm / Kubernetes             & 16\,GB+ & High availability \\
\bottomrule
\end{tabularx}
\end{table}

\section{Complete Installation Procedure}

\subsection{Prerequisites}

\begin{itemize}
  \item Git ;
  \item Docker Engine 24+ and Docker Compose v2.20+ ;
  \item 20\,GB of free disk space (source + Docker images + AI models) ;
  \item Internet access for the first pull of Docker images and Ollama models.
\end{itemize}

\subsection{Installation steps}

\paragraph{Step 1 — Clone the repository}
The source code is publicly available on GitHub:
\begin{lstlisting}[language=YAML]
git clone https://github.com/yousseftaik1/CAIMP.git
cd CAIMP
\end{lstlisting}

\paragraph{Step 2 — Configure environment variables}
\begin{lstlisting}[language=YAML]
cp .env.example .env
# Edit .env: set passwords and generate JWT_SECRET
# JWT_SECRET: openssl rand -hex 64
\end{lstlisting}

\paragraph{Step 3 — Start the platform}
\begin{lstlisting}[language=YAML]
make up        # builds all custom images, runs DB migrations,
               # starts all 19 services (pulls Ollama from Docker Hub)
make ps        # verify every container is "healthy"
\end{lstlisting}

\noindent \texttt{make up} automatically pulls all third-party images (Ollama,
PostgreSQL, Redis, NATS, Grafana, Prometheus, nginx) from Docker Hub on first run.

\paragraph{Step 4 — Download the AI models} (first run only)
\begin{lstlisting}[language=YAML]
make pull-models   # downloads llama3.1:8b-instruct-q4_K_M (~4.7 GB)
                   #          + nomic-embed-text (~270 MB)
\end{lstlisting}

\paragraph{Step 5 — Open the dashboard}
Navigate to \texttt{http://localhost} in a browser.\\
Default credentials: \texttt{admin@caimp.local} / \texttt{Admin1234!}

\paragraph{Step 6 — Enrol a target server}
Generate an enrollment token via the Admin \textsc{api} and run the agent
installer on the target machine (see Section~\ref{sec:vm-monitoring}).

\section{Environment Variables}

\begin{table}[H]
\centering
\caption{Required environment variables}
\label{tab:env-vars}
\small
\begin{tabularx}{\textwidth}{lXl}
\toprule
\textbf{Variable} & \textbf{Description} & \textbf{Required} \\
\midrule
\texttt{POSTGRES\_PASSWORD}    & PostgreSQL password            & Yes \\
\texttt{REDIS\_PASSWORD}       & Redis password                 & Yes \\
\texttt{NATS\_PASSWORD}        & NATS password                  & Yes \\
\texttt{JWT\_SECRET}           & 64-char hex JWT secret         & Yes \\
\texttt{CA\_KEY\_PASSPHRASE}   & Internal CA passphrase         & Yes \\
\texttt{GRAFANA\_ADMIN\_PASSWORD} & Grafana UI password         & Yes \\
\texttt{OLLAMA\_LLM\_MODEL}    & LLM model name                 & No \\
\texttt{SPLUNK\_HOST}          & Splunk host (enables integration) & No \\
\texttt{GITHUB\_WEBHOOK\_SECRET} & GitHub webhook HMAC secret   & No \\
\texttt{LOG\_LEVEL}            & DEBUG/INFO/WARNING             & No \\
\bottomrule
\end{tabularx}
\end{table}

\section{Low-RAM Mode (4 GB)}

The \texttt{docker-compose.test.yml} override applies memory limits on all services,
excludes Ollama by default, replaces Llama 3.1 8B with Llama 3.2 1B (1.3\,GB),
and reduces AI Worker concurrency to 1. A 4\,GB swap file is created automatically
by \texttt{start-test.sh}. Measured RAM usage: \textbf{2.3\,GB} (base stack without Ollama).

\section{Platform Self-Monitoring}

Prometheus scrapes the \texttt{/metrics} endpoint of each service every 15 seconds.
Key metrics monitored: ingestion latency p99, anomaly detection rate,
AI explanation latency, alert deduplication suppression rate, and HTTP 5xx error rate.
All visualised in Grafana at port 3001.

\section{Generating Demo Metrics}

To populate the dashboard with realistic data without needing a real server,
\caimp{} ships a lightweight demo agent (\texttt{services/agent/demo\_agent.py})
that reads actual \textsc{cpu}, \textsc{ram}, and disk metrics from the host
machine and forwards them to the platform every 20 seconds.

\paragraph{Prerequisites}
\begin{lstlisting}[language=YAML]
pip install psutil requests
\end{lstlisting}

\paragraph{Basic usage — one virtual server with real metrics}
\begin{lstlisting}[language=YAML]
python services/agent/demo_agent.py
\end{lstlisting}

\noindent This registers the virtual server \texttt{demo-server-01} in the
default organisation and begins streaming metrics immediately.

\paragraph{Simulating critical alerts — three simultaneous terminals}

\begin{lstlisting}[language=YAML, caption={Terminal 1 — real machine metrics}]
python services/agent/demo_agent.py
\end{lstlisting}

\begin{lstlisting}[language=YAML, caption={Terminal 2 — db-primary with critical RAM}]
CAIMP_SERVER_ID=00000000-0000-0000-0000-000000000013 \
CAIMP_SERVER_NAME=db-primary \
CAIMP_RAM=97 \
python services/agent/demo_agent.py
\end{lstlisting}

\begin{lstlisting}[language=YAML, caption={Terminal 3 — worker-01 with critical CPU}]
CAIMP_SERVER_ID=00000000-0000-0000-0000-000000000014 \
CAIMP_SERVER_NAME=worker-01 \
CAIMP_CPU=99 \
python services/agent/demo_agent.py
\end{lstlisting}

\noindent Within two minutes, the AI Worker generates anomaly explanations for
each critical threshold breach, which appear in the dashboard and are queryable
via the \textsc{ai} chat assistant.

\begin{table}[H]
\centering
\caption{Demo agent environment variable overrides}
\label{tab:demo-agent-vars}
\small
\begin{tabularx}{\textwidth}{lXl}
\toprule
\textbf{Variable} & \textbf{Description} & \textbf{Default} \\
\midrule
\texttt{CAIMP\_SERVER\_ID}   & UUID of the virtual server        & \texttt{...0011} \\
\texttt{CAIMP\_SERVER\_NAME} & Display name in the dashboard     & \texttt{demo-server-01} \\
\texttt{CAIMP\_CPU}          & Override CPU\,\% (0--100)         & real value \\
\texttt{CAIMP\_RAM}          & Override RAM\,\% (0--100)         & real value \\
\texttt{CAIMP\_DISK}         & Override Disk\,\% (0--100)        & real value \\
\texttt{CAIMP\_INTERVAL}     & Sending interval in seconds       & \texttt{20} \\
\texttt{CAIMP\_TW\_URL}      & Telemetry-writer URL              & internal default \\
\bottomrule
\end{tabularx}
\end{table}

\section{Monitoring a Real Virtual Machine}
\label{sec:vm-monitoring}

To monitor a real Linux server (systemd-based), generate an enrollment token via
the Admin \textsc{api}, copy the agent to the target machine, and run the
provided installer.

\paragraph{Step 1 — Generate an enrollment token}
From the machine running \caimp{}:
\begin{lstlisting}[language=YAML]
curl -s -X POST http://localhost/api/v1/admin/servers/enroll \
  -H "Authorization: Bearer <admin_jwt>" \
  -H "Content-Type: application/json" \
  -d '{"name":"my-vm","org_id":"<org_uuid>"}' | python3 -m json.tool
\end{lstlisting}
The response contains a one-time \texttt{enrollment\_token}.

\paragraph{Step 2 — Copy the agent to the target VM}
\begin{lstlisting}[language=YAML]
scp -r services/agent/ user@<vm-ip>:~/caimp-agent/
\end{lstlisting}

\paragraph{Step 3 — Run the installer on the target VM}
\begin{lstlisting}[language=YAML]
ssh user@<vm-ip>
cd ~/caimp-agent
sudo bash install.sh \
  --token <enrollment_token> \
  --url http://<caimp-host-ip>
\end{lstlisting}

\noindent The installer creates a \texttt{caimp-agent} system user, installs a
Python virtual environment, enrolls the agent (generating a cryptographic keypair
and submitting a \textsc{csr}), and starts the \texttt{caimp-agent} systemd
service automatically on boot.

\paragraph{Step 4 — Verify the agent is running}
\begin{lstlisting}[language=YAML]
systemctl status caimp-agent
journalctl -u caimp-agent -f
\end{lstlisting}

\noindent The server appears in the \caimp{} dashboard within 30 seconds of the
first metric batch being received.

\section{Accessing the Admin and Monitoring Dashboards}

\caimp{} runs entirely on the machine where \texttt{docker compose up} was
executed — the administrator's laptop during development, or a dedicated server
in production. All web interfaces are served by the nginx reverse proxy on
port~80.

\begin{table}[H]
\centering
\caption{Web interface access URLs}
\label{tab:web-urls}
\begin{tabularx}{\textwidth}{lXl}
\toprule
\textbf{Interface} & \textbf{Description} & \textbf{URL} \\
\midrule
\caimp{} Dashboard & Main admin portal — login, server list, AI chat
  & \texttt{http://localhost} \\
Grafana & Metrics, time-series graphs, and alert history
  & \texttt{http://localhost:3001} \\
\bottomrule
\end{tabularx}
\end{table}

\paragraph{Network access from another machine}
If \caimp{} runs on a dedicated server with IP \texttt{192.168.1.10}, replace
\texttt{localhost} with that IP address in every URL:
\begin{lstlisting}[language=YAML]
http://192.168.1.10        # CAIMP dashboard
http://192.168.1.10:3001   # Grafana
\end{lstlisting}

\paragraph{Default credentials}
\begin{itemize}
  \item \textbf{\caimp{} dashboard}: \texttt{admin@caimp.local} /
    \texttt{Admin1234!}
  \item \textbf{Grafana}: \texttt{admin} / value of
    \texttt{GRAFANA\_ADMIN\_PASSWORD} in \texttt{.env}
\end{itemize}

\paragraph{Development vs.\ production}
During development the platform runs on the same machine as the browser — no
additional network configuration is required. In production the platform must
run on a server reachable by all monitored machines; the agent installer must
receive the server's public or private IP via the \texttt{--url} flag.

\section*{Conclusion}

The six-step installation procedure makes \caimp{} accessible in under ten minutes,
even on resource-constrained machines. Low-RAM mode and automated backup scripts
ensure operability in constrained environments. The next chapter analyses the
results obtained and discusses future prospects.
